\documentclass{article}
\usepackage[margin=0.65in]{geometry}
\usepackage{enumitem}
\usepackage{graphicx}
\usepackage{hyperref}
\usepackage{titlesec}
\usepackage{xcolor}
\hypersetup{hidelinks}
\setlist[itemize]{leftmargin=1.5em,topsep=2pt,itemsep=1pt,parsep=0pt,partopsep=0pt}
\titlespacing*{\section}{0pt}{1.2ex}{0.6ex}
\setlength{\parindent}{0pt}
\setlength{\parskip}{0.5em}
\setlength{\textfloatsep}{8pt}
\setlength{\intextsep}{6pt}
\setlength{\abovecaptionskip}{4pt}
\setlength{\belowcaptionskip}{0pt}
\makeatletter
\renewcommand{\maketitle}{%
    \begin{center}
        {\Large\bfseries \@title\par}\vspace{0.3em}
        {\normalsize \@author\par}\vspace{0.2em}
        {\small \@date\par}
    \end{center}\vspace{0.4em}
}
\makeatother
% HJNotes-Tags: Basics, Hardware
% HJNotes-Summary: Chip-level MCU vs SoC, and multicore software models SMP/AMP/BMP.
\title{Basic Hardware Concepts}
\author{HJ}
\begin{document}
\maketitle
\section{Chip Level: MCU vs SoC}
\begin{itemize}
    \item \textbf{MCU}: a chip-level computer for deterministic real-time control, with on-chip Flash/SRAM and basic peripherals; typically runs bare-metal firmware or an RTOS, usually has no MMU.
    \item \textbf{SoC}: a system-level chip integrating high-performance processing with complex subsystems (GPU, NPU, DDR, high-speed I/O); typically runs a Rich OS.
\end{itemize}

\textbf{MCU vs.\ SoC}: not an OS-presence distinction. MCU targets deterministic control with on-chip memory; SoC targets complex application processing with external memory and richer integration.

\section{Multicore Software Models: SMP vs AMP vs BMP}
\begin{itemize}
    \item \textbf{SMP (Symmetric)}: one OS manages all cores; tasks may migrate across cores.
    \item \textbf{AMP (Asymmetric)}: separate OS instances run on separate cores/clusters with partitioned resources and inter-core communication.
    \item \textbf{BMP (Bound)}: one OS manages all cores, but tasks/threads are bound to specific cores.
\end{itemize}

\textbf{SMP vs.\ AMP vs.\ BMP}: software usage models for multiple cores, not hardware categories. SMP/BMP share one OS image and differ in task migration; AMP uses multiple OS instances with resource isolation.
\end{document}
